\subsection{$T1$ 定理的陈述与应用}
\label{sec:ch9-t1-statement-applications}

本节陈述主要结果 $T1$ 定理并给出若干应用. 下一节将给出它较长的证明.

考虑把 Schwartz 类映到缓增分布空间的算子
\[
 T:\mathcal S(\mathbb R^n)\longrightarrow\mathcal S'(\mathbb R^n).
\]
换言之, 给定 $f,g\in\mathcal S(\mathbb R^n)$, 已知 $\langle Tf,g\rangle$ 的值. 回忆, 若函数
\[
 K:\mathbb R^n\times\mathbb R^n\setminus\Delta\longrightarrow\mathbb C
\]
满足~\eqref{eq:ch5-standard-kernel-size}、\eqref{eq:ch5-standard-kernel-y-regularity} 和~\eqref{eq:ch5-standard-kernel-x-regularity}, 其中 $\Delta$ 是 $\mathbb R^n\times\mathbb R^n$ 的对角线, 则称 $K$ 为标准核. 若对支集互不相交且紧的 Schwartz 函数 $f,g$ 有
\begin{equation}
 \langle Tf,g\rangle
 =\int_{\mathbb R^n}\int_{\mathbb R^n}
 K(x,y)f(y)g(x)\,dx\,dy,
 \label{eq:ch9-operator-associated-kernel}
\end{equation}
则称 $T$ 与标准核 $K$ 相伴.

伴随算子 $T^*$ 定义为
\[
 \langle T^*f,g\rangle=\langle f,Tg\rangle,
 \qquad f,g\in\mathcal S(\mathbb R^n).
\]
它与标准核 $\widetilde K(x,y)=K(y,x)$ 相伴.

当 $T$ 是 $L^2$ 上的有界算子时, 可把 $f\in L^\infty$ 的 $Tf$ 定义为 BMO 函数, 见定理~\ref{thm:ch6-linfty-to-bmo}. 但这里先验地不知道 $T$ 在 $L^2$ 上有界, 因而不能用该方法定义 $T$ 在 $L^\infty$ 上的作用. 作为替代, 对有界且属于 $C^\infty$ 的函数 $f$ 定义 $Tf$. 事实上, $Tf$ 将是 $C_{c,0}^\infty$ 的对偶元, 其中 $C_{c,0}^\infty$ 是平均为零的 $C_c^\infty$ 函数子空间.

固定 $g\in C_{c,0}^\infty$, 设其支集包含于 $B(0,R)$, 并令 $f\in L^\infty\cap C^\infty$. 取 $\psi_1,\psi_2\in C^\infty(\mathbb R^n)$, 使
\[
 \operatorname{supp}\psi_1\subset B(0,3R),
 \qquad \psi_1=1\text{ on }B(0,2R),
 \qquad \psi_1+\psi_2=1.
\]
则 $f\psi_1\in\mathcal S(\mathbb R^n)$, 所以 $\langle T(f\psi_1),g\rangle$ 有意义. 若 $f$ 具有紧支集, 由~\eqref{eq:ch9-operator-associated-kernel},
\[
 \langle T(f\psi_2),g\rangle
 =\int_{\mathbb R^n}\int_{\mathbb R^n}
 K(x,y)f(y)\psi_2(y)g(x)\,dx\,dy.
\]

\hypertarget{chapter-09-page-214}{}
因为 $f$ 未必具有紧支集, 改为定义
\[
 \langle T(f\psi_2),g\rangle
 =\int_{\mathbb R^n}\int_{\mathbb R^n}
 [K(x,y)-K(0,y)]f(y)\psi_2(y)g(x)\,dx\,dy.
\]
因为 $g$ 的平均为零, 当 $f$ 具有紧支集时这与前一定义一致. 对任意有界 $C^\infty$ 函数 $f$, 该式也有意义, 因为 $y$ 积分远离原点, 且由标准核条件,
\[
 |K(x,y)-K(0,y)|\le\frac{C|x|^\delta}{|y|^{n+\delta}}.
\]
现在定义
\begin{equation}
 \langle Tf,g\rangle
 =\langle T(f\psi_1),g\rangle+\langle T(f\psi_2),g\rangle.
 \label{eq:ch9-t-on-bounded-smooth}
\end{equation}
因为 $g$ 的平均为零, 该定义与 $\psi_1,\psi_2$ 的选择无关. 下文用~\eqref{eq:ch9-t-on-bounded-smooth} 定义 $T1$ 和 $T^*1$.

给定 $f\in C^\infty\cap L^\infty$, 若存在 $b\in\operatorname{BMO}$ 使对所有 $g\in C_{c,0}^\infty$ 有
\[
 \langle Tf,g\rangle=\langle b,g\rangle,
\]
则称 $Tf\in\operatorname{BMO}$. 因为 $C_{c,0}^\infty$ 在 $H^1$ 中稠密, 由 $H^1$ 与 BMO 的对偶性, 即定理~\ref{thm:ch6-h1-dual-bmo}, 这等价于
\[
 |\langle Tf,g\rangle|\le C\|g\|_{H^1}.
\]

为陈述 $T1$ 定理还需一个定义.

\begin{Definition}
\label{def:ch9-weak-boundedness-property}
若对 $C_c^\infty(\mathbb R^n)$ 的每个有界子集 $\mathcal B$, 都存在常数 $C_{\mathcal B}$, 使对任意 $\phi_1,\phi_2\in\mathcal B$, $x\in\mathbb R^n$ 和 $R>0$ 有
\[
 |\langle T\phi_{1,R}^x,\phi_{2,R}^x\rangle|
 \le C_{\mathcal B}R^n,
\]
其中
\[
 \phi_{i,R}^x(y)=\phi_i\left(\frac{y-x}{R}\right),
 \qquad i=1,2,
\]
则称算子 $T$ 具有弱有界性性质, 通常记为 WBP.
\end{Definition}

若算子 $T$ 在某个 $L^p$ 上有界, 则容易证明它具有 WBP.

给定反对称标准核 $K$, 即 $K(x,y)=-K(y,x)$, 可通过
\begin{equation}
 \langle Tf,g\rangle
 =\lim_{\varepsilon\to0}
 \iint_{|x-y|>\varepsilon}K(x,y)f(y)g(x)\,dy\,dx
 \label{eq:ch9-antisymmetric-kernel-operator}
\end{equation}
把它与算子 $T$ 相伴.

\hypertarget{chapter-09-page-215}{}
由 $K$ 的反对称性, 这等价于
\[
 \langle Tf,g\rangle
 =\frac12\iint_{\mathbb R^n\times\mathbb R^n}
 K(x,y)[f(y)g(x)-f(x)g(y)]\,dy\,dx.
\]
该积分绝对收敛. 由中值定理, 在方括号中加上再减去 $f(y)g(y)$, 可得到因子 $|y-x|$, 从而能在对角线附近积分.

由~\eqref{eq:ch9-antisymmetric-kernel-operator} 定义的算子具有 WBP. 只需注意
\[
 \langle T\phi_{1,R}^z,\phi_{2,R}^z\rangle
 =R^n\langle T_{z,R}\phi_1,\phi_2\rangle,
\]
其中 $T_{z,R}$ 按~\eqref{eq:ch9-antisymmetric-kernel-operator} 定义, 核为
\[
 R^nK(Rx+z,Ry+z).
\]
这是一个标准核, 且常数与 $K$ 相同.

一维微分算子说明了 WBP 的重要性. 它把 $\mathcal S(\mathbb R)$ 映到 $\mathcal S'(\mathbb R)$, 与零标准核相伴, 但在 $L^2$ 上无界. 容易验证它不具有 WBP. 微分显然不是通过~\eqref{eq:ch9-antisymmetric-kernel-operator} 与零核相伴的算子, 后者是平凡地具有 WBP 的零算子.

\begin{Theorem}
\label{thm:ch9-t1}
设
\[
 T:\mathcal S(\mathbb R^n)\longrightarrow\mathcal S'(\mathbb R^n)
\]
与标准核 $K$ 相伴. 则 $T$ 可延拓为 $L^2(\mathbb R^n)$ 上的有界算子当且仅当:
\begin{enumerate}[label=(\arabic*)]
\item $T1\in\operatorname{BMO}$;
\item $T^*1\in\operatorname{BMO}$;
\item $T$ 具有 WBP.
\end{enumerate}
\end{Theorem}

这三个条件的必要性已经得到. (1) 和 (2) 来自定理~\ref{thm:ch6-linfty-to-bmo} 及其后的注记, (3) 的必要性在定义~\ref{def:ch9-weak-boundedness-property} 后已说明.

\begin{Corollary}
\label{cor:ch9-antisymmetric-t1}
若 $K$ 是反对称标准核, $T$ 是由~\eqref{eq:ch9-antisymmetric-kernel-operator} 定义的算子, 则 $T$ 在 $L^2(\mathbb R^n)$ 上有界当且仅当 $T1\in\operatorname{BMO}$.
\end{Corollary}

这立即由定理~\ref{thm:ch9-t1} 得到, 因为 $T$ 总具有 WBP, 且 $T^*1=-T1$.

\begin{Example}[An application of the $T1$ theorem]
\label{ex:ch9-calderon-commutators}
回忆第~\ref{sec:ch5-generalized-calderon-zygmund}~节定义的 Calderón 交换子
\[
 T_kf(x)=\lim_{\varepsilon\to0}
 \int_{|x-y|>\varepsilon}
 \left(\frac{A(x)-A(y)}{x-y}\right)^k\frac{f(y)}{x-y}\,dy,
\]
其中 $A:\mathbb R\to\mathbb R$ 是 Lipschitz 函数, $k\ge0$ 是整数. 这些算子由~\eqref{eq:ch9-antisymmetric-kernel-operator} 定义, 其反对称核为
\[
 K_k(x,y)=\frac{(A(x)-A(y))^k}{(x-y)^{k+1}}.
\]
\end{Example}

\hypertarget{chapter-09-page-216}{}
\begin{Corollary}
\label{cor:ch9-calderon-commutator-bounds}
算子 $T_k$ 在 $L^2$ 上有界, 且存在正常数 $C$ 使
\[
 \|T_k\|\le C^k\|A'\|_\infty^k.
\]
\end{Corollary}

由该结果和定义~\ref{def:ch5-calderon-zygmund-operator}, $T_k$ 是 Calderón--Zygmund 算子.

\begin{Proof}
由定理~\ref{thm:ch9-t1} 的证明将看到, 算子 $T$ 的范数线性依赖于假设中涉及的常数. 在推论~\ref{cor:ch9-antisymmetric-t1} 的特殊情形中, 这些常数归结为 $K$ 作为标准核的常数以及 $T_k1$ 的 BMO 范数.

直接论证表明 $K_k$ 作为标准核的常数由 $C_1(k+1)\|A'\|_\infty^k$ 控制. 所以只需证明对某个 $C>0$ 和所有 $k$,
\begin{equation}
 \|T_k1\|_*\le C^{k+1}\|A'\|_\infty^k.
 \label{eq:ch9-commutator-bmo-induction}
\end{equation}
若 $C\ge C_1$, 则
\[
 \|T_k\|\le C_2(C_1(k+1)+C^{k+1})\|A'\|_\infty^k
 \le2C_2C^{k+1}\|A'\|_\infty^k.
\]

对~\eqref{eq:ch9-commutator-bmo-induction} 作归纳证明. 当 $k=0$ 时, $T_0$ 是 Hilbert 变换, 因而 $T_01=0$. 设~\eqref{eq:ch9-commutator-bmo-induction} 对某个 $k$ 成立. 分部积分给出
\[
 T_{k+1}1=T_kA'.
\]
形式上该计算是直接的, 但严格化需要一些细节, 留给读者. 因此
\[
 \|T_{k+1}1\|_*=\|T_kA'\|_*
 \le\|T_k\|_{L^\infty,\operatorname{BMO}}\|A'\|_\infty.
\]
由定理~\ref{thm:ch6-linfty-to-bmo} 的证明,
\[
 \begin{aligned}
 \|T_k\|_{L^\infty,\operatorname{BMO}}
 &\le C_3(C_1(k+1)\|A'\|_\infty^k+\|T_k\|_{L^2,L^2})\\
 &\le C_3(C_1(k+1)+2C_2C^{k+1})\|A'\|_\infty^k.
 \end{aligned}
\]
因为 $C_1,C_2,C_3$ 均与 $k$ 无关, 当 $C$ 足够大时
\[
 C_3(C_1(k+1)+2C_2C^{k+1})\le C^{k+2}.
\]
这给出所需不等式.
\end{Proof}

\hypertarget{chapter-09-page-217}{}
第~\ref{sec:ch5-generalized-calderon-zygmund}~节还考虑了与反对称核
\[
 K(x,y)=\frac1{x-y+i(A(x)-A(y))}
\]
相伴的另一个算子, 它与 Lipschitz 曲线上的 Cauchy 积分有关. 已经证明, 若 $\|A'\|_\infty<1$, 则可把该核展开为
\[
 K(x,y)=\sum_{k=0}^\infty i^kK_k(x,y).
\]
由该展开和推论~\ref{cor:ch9-calderon-commutator-bounds} 得到:

\begin{Corollary}
\label{cor:ch9-small-lipschitz-cauchy}
存在 $\varepsilon>0$, 使得若 $\|A'\|_\infty\le\varepsilon$, 则与核 $K$ 相伴的算子是 Calderón--Zygmund 算子.
\end{Corollary}

关于这些算子的进一步内容见第~\ref{subsec:ch9-calderon-cauchy}~节.
